\section{问题一：多模态特征提取与时间统一}

\subsection{问题定义与数据结构}
附件 1 包含 100 条视频样本。每条样本记录视频编号、文本转写、原始视频时长以及文本、音频和视觉信息。本文不将问题一表述为预测性能问题，而是建立从原始样本到统一多模态时序张量的可复现映射。

\subsection{三模态特征}
文本使用预训练语言模型提取上下文表示，音频提取 MFCC、Chroma、Mel 频谱和相关谱统计特征，视觉提取视频帧的灰度直方图及统计特征。每条记录同时保存特征维度、原始有效长度和异常处理状态。问题一生成的特征文件和 100 条样本汇总表用于后续问题二、问题三。

\subsection{统一长度与时间位置}
设文本、音频和视觉的原始序列分别为 $X_t,X_a,X_v$。本文将文本有效序列作为统一索引基准，将音频和视觉按时间顺序划分为对应区间并进行平均聚合：
\begin{equation}
\widetilde{x}_{m,k}=\frac{1}{|I_k|}\sum_{i\in I_k}x_{m,i},\qquad m\in\{a,v\},
\end{equation}
其中 $I_k$ 表示映射到第 $k$ 个统一位置的原始特征区间。统一长度设为 $L=50$；超过长度的序列截断，不足长度的位置零填充，并保存有效位置掩码 $M_m(k)$。

文本掩码来自原始有效 token；音频和视觉掩码来自特征提取阶段记录的有效长度或缺失标记。这样可以区分真实缺失、填充位置和数值接近零但有效的特征，避免仅凭特征绝对值推断有效性。

\subsection{质量检查与输出}
对齐后逐条检查样本编号、特征维度、序列长度、有效位置数量和标签数量。异常视频、空音频和截断情况写入处理日志。最终输出包括 100 条样本汇总表、统一特征文件、有效长度和掩码字段。

\subsection{问题一小结}
问题一的成果是可复现的数据转换流程，而不是单独的分类指标。该流程保证文本、音频和视觉特征具有统一的批处理形状，并为后续缺失模态实验提供明确的 mask 和有效长度定义。
